%=============================================================
%  Chapter 7 --- Iteration 3 : Administration & Security
%=============================================================
\chapter{Implementation --- Iteration 3: Administration and Security}

\section*{Introduction}
The third and final iteration consolidates the platform on two fronts:
a complete \textbf{administration dashboard} enabling operators to manage
users, orders, and platform statistics, and a \textbf{security hardening layer}
covering JWT token revocation, role-based access control, and strict input
validation across all endpoints.
At the end of this iteration, \textbf{LUMINA} is ready for production deployment.

%-------------------------------------------------------------
\section{Sprint Planning}
%-------------------------------------------------------------

\subsection{Sprint Goal}

\begin{tcolorbox}[colback=blue!5, colframe=blue!40, title=Sprint Goal --- Iteration 3]
Give administrators the tools they need to operate the platform daily, while
ensuring that no unauthorized action can be performed by a malicious user
or a compromised token.
\end{tcolorbox}

\subsection{Sprint Backlog}

\begin{table}[H]
\centering
\caption{User stories --- Iteration 3}
\begin{tabularx}{\textwidth}{|c|X|c|c|}
\hline
\rowcolor{gray!20}
\textbf{ID} & \textbf{User Story} & \textbf{Priority} & \textbf{Points} \\
\hline
US-17 & As an admin, I want a dashboard showing key metrics (orders, revenue, users). & Must & 8 \\
\hline
US-18 & As an admin, I want to list, search, and view the detailed profile of each user. & Must & 5 \\
\hline
US-19 & As an admin, I want to suspend or reactivate a user account. & Must & 5 \\
\hline
US-20 & As an admin, I want to list all orders and update their status. & Must & 5 \\
\hline
US-21 & As a customer, I want my session immediately invalidated on logout, even if someone holds my token. & Must & 8 \\
\hline
US-22 & As a system, all protected routes must verify the caller's role and reject revoked tokens. & Must & 8 \\
\hline
US-23 & As a system, all incoming data must be validated and typed before reaching the business layer. & Must & 5 \\
\hline
US-24 & As a super admin, I want to promote a user to the admin role. & Should & 3 \\
\hline
\end{tabularx}
\end{table}

%-------------------------------------------------------------
\section{Design}
%-------------------------------------------------------------

\subsection{Role Model}

The platform distinguishes three access levels:

\begin{table}[H]
\centering
\caption{Role and permissions matrix}
\begin{tabularx}{\textwidth}{|X|c|c|c|}
\hline
\rowcolor{gray!20}
\textbf{Permission} & \textbf{Customer} & \textbf{Admin} & \textbf{Super Admin} \\
\hline
Catalogue, cart, orders & \checkmark & \checkmark & \checkmark \\
\hline
Product management & & \checkmark & \checkmark \\
\hline
User management & & \checkmark & \checkmark \\
\hline
Order management & & \checkmark & \checkmark \\
\hline
Statistics dashboard & & \checkmark & \checkmark \\
\hline
Promote to admin role & & & \checkmark \\
\hline
Platform configuration & & & \checkmark \\
\hline
\end{tabularx}
\end{table}

\subsection{Security Architecture}

Security is implemented in successive layers. Every request passes through
the following filters in order:

\begin{enumerate}
  \item \textbf{HTTPS} --- transport encryption (TLS).
  \item \textbf{JWT validation} --- signature, expiration, \texttt{jti} claim.
  \item \textbf{Redis blacklist} --- reject if the \texttt{jti} is revoked.
  \item \textbf{Role check} --- verify the \texttt{role} field in the payload.
  \item \textbf{Input validation (Zod)} --- typed DTOs before reaching the use case.
  \item \textbf{Business logic} --- isolated use cases with no direct DB access.
\end{enumerate}

%-------------------------------------------------------------
\section{Implementation}
%-------------------------------------------------------------

\subsection{Admin Dashboard}

The \texttt{/admin} page displays key platform indicators:
today's orders, monthly revenue, new users, and low-stock products.
Data is aggregated by direct PostgreSQL queries from the
\texttt{AdminController}, without a caching layer to ensure freshness.

\begin{figure}[H]
  \centering
  \includegraphics[width=\textwidth]{figures/screen_admin_dashboard.png}
  \caption{Admin dashboard --- global statistics}
  \label{fig:screen_admin_dashboard}
\end{figure}

%-------------------------------------------------------------
\subsection{User Management}
%-------------------------------------------------------------

\subsubsection{List and Search}

The \texttt{/admin/users} interface lists all accounts with pagination,
role filtering, and search by name or email.
Each row shows account status (active / suspended), registration date,
and order count.

\begin{figure}[H]
  \centering
  \includegraphics[width=\textwidth]{figures/screen_admin_users.png}
  \caption{Admin interface --- user list}
  \label{fig:screen_admin_users}
\end{figure}

\subsubsection{User Profile}

A user's detail page (\texttt{/admin/users/[id]}) is loaded via a dedicated
\texttt{GET /api/admin/users/:userId} endpoint, avoiding fetching the full
list client-side just to find one entry.

\begin{lstlisting}[language=Java, caption=User detail endpoint (AdminController)]
@Get('/users/:userId')
@Authorized(['admin', 'super_admin'])
async getUserById(@Param('userId') userId: string, @Res() res: Response) {
  const user = await db.select().from(userTable)
    .where(eq(userTable.id, userId)).limit(1);
  if (!user.length)
    return res.status(404).json({ message: 'User not found' });
  return res.json(user[0]);
}
\end{lstlisting}

\begin{figure}[H]
  \centering
  \includegraphics[width=0.85\textwidth]{figures/screen_admin_user_detail.png}
  \caption{Admin interface --- user detail page}
  \label{fig:screen_admin_user_detail}
\end{figure}

\subsubsection{Suspend and Reactivate}

Suspending an account sets the \texttt{isSuspended = true} flag in the
database. The authentication middleware checks this flag on every request:
a suspended user receives a \texttt{403 Forbidden} response even with a
valid token.

\begin{lstlisting}[language=Java, caption=SuspendUserUseCase]
async execute({ targetUserId, suspend }: SuspendUserInput): Promise<Result<void>> {
  const user = await this.userRepo.findById(targetUserId);
  if (!user) return ResultHelper.failure('User not found', ErrorCode.NOT_FOUND);
  if (user.role === 'super_admin')
    return ResultHelper.failure('Cannot suspend a super admin', ErrorCode.FORBIDDEN);
  user.isSuspended = suspend;
  await this.userRepo.update(user);
  return ResultHelper.success(undefined);
}
\end{lstlisting}

%-------------------------------------------------------------
\subsection{Order Management}
%-------------------------------------------------------------

The \texttt{/admin/orders} interface lists all orders with status filtering.
A dropdown lets the admin advance an order's status
(\textsc{confirmed} -> \textsc{shipped} -> \textsc{delivered})
or cancel it (\textsc{cancelled}).

\begin{figure}[H]
  \centering
  \includegraphics[width=\textwidth]{figures/screen_admin_orders.png}
  \caption{Admin interface --- order management}
  \label{fig:screen_admin_orders}
\end{figure}

%-------------------------------------------------------------
\subsection{Security --- JWT Blacklist via Redis}
%-------------------------------------------------------------

On logout, the user's JWT is immediately invalidated, even if it has not
yet expired. The mechanism relies on the \texttt{jti} (JWT ID) claim ---
a UUID v4 unique to each token, included at issuance.

\begin{lstlisting}[language=Java, caption=Token revocation on logout]
// blacklistToken(rawToken)
const decoded = JwtService.staticDecode(rawToken) as { jti: string; exp: number };
const ttl = decoded.exp - Math.floor(Date.now() / 1000);
if (ttl > 0)
  await redisClient.set(`blacklist:jti:${decoded.jti}`, '1', { EX: ttl });
\end{lstlisting}

\begin{lstlisting}[language=Java, caption=Blacklist check on every protected request]
// tokenBlacklistMiddleware --- fail CLOSED
const revoked = await redisClient.get(`blacklist:jti:${payload.jti}`);
if (revoked !== null)
  return res.status(401).json({ error: 'Token has been revoked' });
// Redis unavailable -> 503, never 200
\end{lstlisting}

The \textit{fail-closed} behavior ensures that if Redis is unavailable,
no potentially revoked token is accepted --- the service returns
\texttt{503 Service Unavailable} rather than letting the request through.

%-------------------------------------------------------------
\subsection{Security --- Role-Based Access Control (RBAC)}
%-------------------------------------------------------------

All administration routes are protected by the \texttt{@Authorized} decorator
from \textit{routing-controllers}, which verifies the role extracted from the
JWT payload after blacklist validation.

\begin{lstlisting}[language=Java, caption=Admin routes protected by role]
@JsonController('/admin')
@UseBefore(tokenBlacklistMiddleware)
export class AdminController {

  @Get('/users')
  @Authorized(['admin', 'super_admin'])
  async listUsers(@Res() res: Response) { ... }

  @Patch('/users/:id/role')
  @Authorized(['super_admin'])           // super admin only
  async assignRole(...) { ... }
}
\end{lstlisting}

%-------------------------------------------------------------
\subsection{Security --- Input Validation (Zod DTOs)}
%-------------------------------------------------------------

Every endpoint accepting input data is associated with a Zod-validated DTO
before reaching the corresponding use case.
Validation errors are normalized into a
\texttt{422 Unprocessable Entity} response with per-field details.

\begin{lstlisting}[language=Java, caption=Example DTO --- CheckoutDto.ts]
export const checkoutSchema = z.object({
  shippingAddress: z.string().min(10, 'Address too short'),
  paymentMethod:   z.enum(['cash_on_delivery']),
  currency:        z.string().length(3).default('TND'),
});
export type CheckoutDto = z.infer<typeof checkoutSchema>;
\end{lstlisting}

DTOs cover all input endpoints: product creation/editing, cart operations,
order placement, and admin actions.

%-------------------------------------------------------------
\subsection{Role Promotion (Super Admin)}
%-------------------------------------------------------------

A super administrator can promote a user to the admin role from the user
management interface. This action is restricted to the \texttt{super\_admin}
role and logged server-side for auditability.

\begin{figure}[H]
  \centering
  \includegraphics[width=0.75\textwidth]{figures/screen_admin_role.png}
  \caption{Super admin interface --- role promotion}
  \label{fig:screen_admin_role}
\end{figure}

%-------------------------------------------------------------
\section{Sprint Review}
%-------------------------------------------------------------

\subsection{Delivered Features}

\begin{table}[H]
\centering
\caption{Iteration 3 results}
\begin{tabular}{|c|l|c|}
\hline
\rowcolor{gray!20}
\textbf{ID} & \textbf{User Story} & \textbf{Status} \\
\hline
US-17 & Statistics dashboard & \textcolor{green!60!black}{\checkmark\ Done} \\
\hline
US-18 & User list and search & \textcolor{green!60!black}{\checkmark\ Done} \\
\hline
US-19 & Account suspension / reactivation & \textcolor{green!60!black}{\checkmark\ Done} \\
\hline
US-20 & Order management & \textcolor{green!60!black}{\checkmark\ Done} \\
\hline
US-21 & JWT revocation on logout & \textcolor{green!60!black}{\checkmark\ Done} \\
\hline
US-22 & RBAC + blacklist check on all routes & \textcolor{green!60!black}{\checkmark\ Done} \\
\hline
US-23 & Zod validation on all endpoints & \textcolor{green!60!black}{\checkmark\ Done} \\
\hline
US-24 & Role promotion (super admin) & \textcolor{green!60!black}{\checkmark\ Done} \\
\hline
\end{tabular}
\end{table}

\subsection{Challenges}

\begin{itemize}
  \item \textbf{Redis fail-closed behavior}: the initial temptation was to
        pass requests through when Redis was unavailable to avoid degrading
        the user experience. This approach was rejected: a revoked token
        (e.g., after session theft) must \textit{always} be blocked;
        availability is sacrificed in favor of security.
  \item \textbf{Redis TTL consistency}: the blacklist TTL must be calculated
        relative to the token's expiration, not a fixed duration. A small clock
        skew between the token issuer and Redis required adding a safety margin.
  \item \textbf{Pagination}: the user and order lists can become large.
        Drizzle offset/limit pagination was preferred over cursor-based
        pagination for implementation simplicity at this stage.
\end{itemize}

\section*{Conclusion}
Iteration 3 completed the construction of \textbf{LUMINA} by delivering the
administration tools needed for daily platform operation, alongside a
multi-layered security architecture covering authentication, session revocation,
access control, and data validation.
All three iterations collectively covered the full product backlog defined at
the start of the project, delivering a secure, intelligent, and production-ready
e-commerce platform.
